% ############################################################################
\section{Containerized development workflows}

% Related documentation files:
% - helpers_root/docs/tools/docker/all.docker.tutorial.md
%   Comprehensive Docker tutorial covering development workflows
% - helpers_root/docs/tools/docker/all.dockerized_flow.explanation.md
%   Explains the containerized development flow
% - helpers_root/docs/tools/docker/all.dind_and_sibling_containers.how_to_guide.md
%   Details on Docker-in-Docker and sibling container patterns
% - helpers_root/docs/tools/docker/all.develop_in_helpers.how_to_guide.md
%   Developing within containerized environments
% - helpers_root/docs/tools/thin_environment/all.thin_environment.reference.md
%   Reference for the thin environment concept
% - helpers_root/docs/tools/thin_environment/all.gh_thin_env_dependencies.how_to_guide.md
%   Managing thin environment dependencies
% - helpers_root/docs/tools/thin_environment/all.gh_and_thin_env_requirements.reference.md
%   Requirements and specifications for thin environments
% - helpers_root/docs/tools/dev_system/all.dealing_with_missing_dependencies.how_to_guide.md
%   Troubleshooting container dependency issues

A ``thin environment'' contains the minimal dependency setup to bootstrap
development and deployment using ``runnable directories" 
across diverse systems (servers, laptops, CI/CD). Runnable directories are a
Docker-based system that ensures reproducible environments across development,
testing, and production stages, with support for both children and sibling
container patterns.

% TODO(gp):

% ############################################################################
\section{The \texttt{//helpers} submodule}

% Related documentation files:
% - helpers_root/docs/tools/dev_system/all.create_a_super_repo_with_helpers.how_to_guide.md
%   Explains how to create a repository that includes helpers as submodule
% - helpers_root/docs/tools/dev_system/all.integrate_repos.how_to_guide.md
%   Guide on integrating multiple repositories using helpers
% - helpers_root/docs/tools/dev_system/all.replace_common_files_with_script_links.md
%   Details the symbolic link strategy for sharing common files
% - helpers_root/docs/tools/git/all.integrate_helpers.how_to_guide.md
%   Instructions for integrating helpers submodule into repositories
% - helpers_root/docs/onboarding/all.development_setup.how_to_guide.md
%   Initial setup including helpers submodule installation

Helpers is a foundational Python utility library that serves as the core
infrastructure for an ecosystem of software development repositoriesIt's
designed to be integrated as a git submodule into ``super-repos" that extend
its functionality.

It contains all the code for all the functionalities described in this
paper.

% ############################################################################
\section{Runnable directory architecture}

% Related documentation files:
% - helpers_root/docs/tools/dev_system/all.runnable_repo.reference.md
%   Core reference for runnable directory architecture
% - helpers_root/docs/tools/dev_system/all.create_a_runnable_dir.how_to_guide.md
%   Step-by-step guide to creating a runnable directory
% - helpers_root/docs/tools/dev_system/all.devops_docker.reference.md
%   Reference for DevOps setup in runnable directories
% - helpers_root/docs/tools/dev_system/all.devops_docker.how_to_guide.md
%   How to set up DevOps workflows for runnable directories
% - docs/kaizenflow/all.install_helpers.how_to_guide.md
%   Installing and using helpers across multiple runnable directories
% - infra/docs/operations/ck.infra_runnable_dir.how_to_guide.md
%   Infrastructure setup for runnable directories

% TODO(gp): Merge in the full paper

% ============================================================================
\subsection{Problem}

Software development teams face a fundamental organizational dilemma when
structuring their codebases. The choice between monorepos and multi-repos
presents a difficult trade-off with no universally satisfactory solution.

Monorepos provide environment consistency and simplified version control,
but struggle with scalability.

Multi-repos offer modularity and independent development lifecycles, but
introduce synchronization issues between repositories.

The core challenge is achieving both modularity and consistency
simultaneously---allowing teams to work independently on separate components
while ensuring all components operate in compatible, reproducible environments.
Traditional approaches force teams to sacrifice one for the other.

% ============================================================================
\subsubsection{The monorepo approach}

% TODO(gp): Fix the references

The monorepo approach involves storing all code for multiple applications
within a single repository. This strategy has been popularized by
large tech companies like Google[2], Meta[3], Microsoft[4] and Uber[5],
proving that even codebases with billions of lines of code can be effectively
managed in a single repository. The key benefits of this approach
include:

\begin{itemize}
  \item Consistency in environment: with everything housed in one
    repository, there's no risk of projects becoming incompatible due to
    conflicting versions of third-party packages.

  \item Simplified version control: there is a single commit history, which
    makes it easy to track and, if needed, revert changes globally.

  \item Reduced coordination overhead: developers work within the same
    repository, with easy access to all code, shared knowledge, tools and
    consistent coding standards.
\end{itemize}

However, as monorepo setups scale, users often face significant challenges.
A major downside is long CI/CD build times, as even small changes can trigger
massive rebuilds and tests throughout the entire codebase. To cope with this,
extra tooling, such as [Buck](https://buck2.build/) or
[Bazel](https://bazel.build/), must be configured, adding complexity to
workflows. Even something as simple as searching and browsing the code becomes
more difficult, often requiring specialized tools and IDE plug-ins.

Additionally, when everything is located in one place, it is harder to
separate concerns and maintain clear boundaries between projects. Managing
permissions also becomes more difficult when only selected developers should
have access to specific parts of the codebase.

% ============================================================================
\subsubsection{The monorepo approach}

The multi-repo approach involves splitting code across several
repositories, with each one dedicated to a specific module or service.
This modularity allows teams to work independently on different parts
of a system, making it easier to manage changes and releases for individual
components. Each repository can evolve at its own pace, and developers
can focus on smaller, more manageable codebases.

However, the multi-repo strategy comes with its own set of challenges,
particularly when it comes to managing dependencies and ensuring
version compatibility across repositories. For instance, different
repositories might rely on two different versions of a third-party
package, or even conflicting packages, making synchronization complex
or, in some cases, nearly impossible. In general, propagating changes
from one repository to another requires careful coordination. Tools like
[Jenkins](https://www.jenkins.io/) and [GitHub Actions](https://github.com/features/actions)
help streamline CI/CD pipelines, but they often struggle when dealing with
heterogeneous environments.

% ============================================================================
\subsection{Solution: Runnable directories}

An ideal strategy would combine the best of both worlds:

\begin{itemize}
  \item The modularity of multi-repos, to keep the codebase scalable
    and simplify day-to-day development processes.

  \item The environment consistency of monorepos, to avoid
    synchronization issues and prevent errors that arise from
    executing code in misaligned environments.
\end{itemize}

Both are achieved through the hybrid approach proposed in this paper,
which will be discussed in Section 3.

\begin{itemize}
  \item This section describes the design principles in our approach to
    create Git repos that contain code that can be:

    \begin{itemize}
      \item Composed through different Git sub-module

      \item Tested, built, run, and released (on a per-directory basis
        or not)
    \end{itemize}

  \item The technologies that this approach relies on are:

    \begin{itemize}
      \item Git for source control

      \item Python virtual environment and \texttt{poetry} (or similar)
        to control Python packages

      \item \texttt{pytest} for unit and end-to-end testing

      \item Docker for managing containers
    \end{itemize}

  \item The approach described in this paper is not strictly dependent
    of the specific package (e.g., \texttt{poetry} can be replaced by \texttt{Conda}
    or another package manager)
\end{itemize}

% ============================================================================
\subsection{Design goals}

The proposed development system supports the following functionalities

% ----------------------------------------------------------------------------
\subsubsection{Development functionalities}

\begin{itemize}
  \item Support composing code using a GitHub sub-module approach

  \item Make it easy to add the development tool chain to a ``new project''
    by simply adding the Git sub-module \texttt{//helpers} to the
    project

  \item Create complex workflows (e.g., for dev and devops
    functionalities) using makefile-like tools based on Python \texttt{invoke}
    package

  \item Have a simple way to maintain common files across different repos
    in sync through links and automatically diff-ing files

  \item Support for both local and remote development using IDEs (e.g.,
    PyCharm, Visual Studio Code)
\end{itemize}

% ----------------------------------------------------------------------------
\subsubsection{Python package management}

\begin{itemize}
  \item Carefully manage and control dependencies using Python managers
    (such as \texttt{poetry}) and virtual environments

  \item Code and containers can be versioned and kept in sync automatically
    since a certain version of the code can require a certain version of
    the container to run properly

    \begin{itemize}
      \item Code is versioned through Git

      \item Each container has a \texttt{changelog.txt} that contains
        the current version and the history
    \end{itemize}
\end{itemize}

% ----------------------------------------------------------------------------
\subsubsection{Testing}

\begin{itemize}
  \item Run end-to-end tests using \texttt{pytest} by automatically
    discover tests based on dependencies and test lists, supporting the
    dependencies needed by different directories

  \item Native support for both children-containers (i.e., Docker-in-Docker)
    and sibling containers
\end{itemize}

% ----------------------------------------------------------------------------
\subsubsection{DevOps functionalities}

\begin{itemize}
  \item Support automatically different stages for container development

    \begin{itemize}
      \item E.g., \texttt{test} / \texttt{local}, \texttt{dev},
        \texttt{prod}
    \end{itemize}

  \item Standardize ways of building, testing, retrieving, and
    deploying containers

  \item Ensure alignment between development environment, deployment, and
    CI/CD systems (e.g., GitHub Actions)

  \item Bootstrap the development system using a ``thin environment'',
    which has the minimum number of dependencies to allow development
    and deployment in exactly the same way in different setups (e.g., server,
    personal laptop, CI/CD)

  \item Manage dependencies in a way that is uniform across platforms and
    OSes, using Docker containers

  \item Separate the need to:

    \begin{itemize}
      \item Build and deploy containers (by devops)

      \item Use containers to develop and test (by developers)
    \end{itemize}

  \item Built-in support for multi-architecture builds (e.g, for Intel
    \texttt{x86} and Arm) across different OSes supporting containers (e.g.,
    Linux, MacOS, Windows Subsystem for Linux WSL)

  \item Support for developing, testing, and deploying multi-container
    applications
\end{itemize}

% ----------------------------------------------------------------------------
\subsubsection{Runnable directory}

The core concept of the proposed approach is a \textbf{runnable
directory} --- a self-contained, independently executable directory with
code, equipped with a dedicated DevOps setup. A repository is thus a special
case of a runnable directory. Developers typically work within a
single runnable directory for a given application, enabling them to test
and deploy code without affecting other parts of the codebase.

A runnable directory can contain other runnable directories as subdirectories.
For example, Figure 1 depicts three runnable directories: \texttt{A},
\texttt{B}, and \texttt{C}. Here, \texttt{A} and \texttt{C} are repositories,
with \texttt{C} incorporated into \texttt{A} as a submodule, while \texttt{B}
is a subdirectory within \texttt{A}. This setup provides the same
accessibility as if all the code were hosted in a single monorepo. Note
that each of \texttt{A}, \texttt{B}, and \texttt{C} has its own DevOps
pipeline --- a key feature of our approach, which is discussed further
in Section 3.2.

% rendered_images:begin
% ```mermaid
% graph RL
%     subgraph A [Runnable Dir A]
%         direction TB
%         subgraph C1 [Runnable Dir C]
%             DevOpsC1[DevOps C]
%         end
%         subgraph B [Runnable Dir B]
%             DevOpsB[DevOps B]
%         end
%         DevOpsA[DevOps A]
%     end
%     subgraph C [Runnable Dir C]
%         DevOpsC[DevOps C]
%     end
% 
%     C -->|Submodule| C1
% 
%     style A fill:#FFF3CD
%     style C fill:#FFF3CD,stroke:#9E9D24
% ```
% label=fig:Runnable_Dir_Arch
% caption=Sample architecture of Causify's runnable directories.
% rendered_images:end
% render_images:begin
\begin{figure}[H]
  \includegraphics[width=\linewidth]{part2.runnable_directory.tex.figs/part2.runnable_directory.1.png}
  \caption{Sample architecture of Causify's runnable directories.}
  \label{fig:Runnable_Dir_Arch}
\end{figure}
% render_images:end

% ----------------------------------------------------------------------------
\subsubsection{Docker}

Docker is the backbone of our containerized development environment. Every
runnable directory contains Dockerfiles that allow it to build and run
its own Docker containers, which include the code, its dependencies,
and the runtime system.

This Docker-based approach addresses two important challenges. First,
it ensures consistency by isolating the application from variations in
the host operating system or underlying infrastructure. Second, a
specific package (or package version) can be added to the container of
a particular runnable directory without affecting other parts of the
codebase. This prevents ``bloating'' the environment with packages required
by all applications --- a common issue in monorepos --- while also effectively
mitigating the risk of conflicting dependencies, which can arise in a
multi-repo setup.

Our approach supports multiple stages for container release:

\begin{itemize}
  \item Local: used to work on updates to the container; accessible only
    to the developer who built it.

  \item Development: used by all team members in day-to-day
    development of new features.

  \item Production: used to run the system by end users.
\end{itemize}

This multi-stage workflow enables seamless progression from testing to
system deployment.

It is also possible to run a container within another container's environment
in a Docker-in-Docker setup. In this case, children containers are started
directly inside a parent container, allowing nested workflows or
builds. Alternatively, sibling containers can run side by side and
share resources such as the host's Docker daemon, enabling inter-container
communication and orchestration.

% rendered_images:begin
% ```mermaid
% graph TD
%     host[Host]
%     docker_engine[Docker Engine]
%     subgraph sibling_container["Sibling Containers"]
%         container_1[Container 1]
%         container_2[Container 2]
%     end
%     subgraph children_container["Children Containers"]
%         container_1a[Container 1a]
%         container_1b[Container 1b]
%     end
%     host --> docker_engine
%     docker_engine --> container_1
%     docker_engine --> container_2
%     container_1 --> container_1a
%     container_1 --> container_1b
% 
%     style sibling_container fill:#FFF3CD,stroke:#9E9D24
%     style children_container fill:#FFF3CD,stroke:#9E9D24
% ```
% caption=Docker container flow.
% label=fig:Docker_Flow
% rendered_images:end
% render_images:begin
\begin{figure}[H]
  \includegraphics[width=\linewidth]{part2.runnable_directory.tex.figs/part2.runnable_directory.2.png}
  \caption{Docker container flow.}
  \label{fig:Docker_Flow}
\end{figure}
% render_images:end

% ----------------------------------------------------------------------------
\subsubsection{What is needed}

An ideal strategy would combine the best of both worlds:

\begin{itemize}
  \item The modularity of multi-repos, to keep the codebase scalable
    and simplify day-to-day development processes.

  \item The environment consistency of monorepos, to avoid
    synchronization issues and prevent errors that arise from
    executing code in misaligned environments.
\end{itemize}

Both are achieved through the hybrid approach proposed in this paper,
which will be discussed in Section 3.

% ----------------------------------------------------------------------------
\subsubsection{Thin environment}

To bootstrap development workflows, we use a thin client that installs
a minimal set of essential dependencies, such as Docker and invoke, in
a lightweight virtual environment. A single thin environment is shared
across all runnable directories which minimizes setup overhead (see Figure
3). This environment contains everything that is needed to start development
containers, which are in turn specific to each runnable directory. With
this approach, we ensure that development and deployment remain
consistent across different systems (e.g., server, personal laptop, CI/CD).

% rendered_images:begin
% ```mermaid
% graph RL
%   thin_env[thin environment]
%   subgraph A [Runnable Dir A]
%     direction TB
%         B[Runnable Dir B]
%         C1[Runnable Dir C]
%     end
%     subgraph C [Runnable Dir C]
%     end
% 
%   C -->|Submodule| C1
%   A -.-> thin_env
%   B -.-> thin_env
%   C1 -.-> thin_env
%   C -.-> thin_env
% 
%   style A fill:#FFF3CD
%   style C fill:#FFF3CD,stroke:#9E9D24
% ```
% caption=Thin environment shared across multiple runnable directories.
% rendered_images:end
% render_images:begin
\begin{figure}[H]
  \includegraphics[width=\linewidth]{part2.runnable_directory.tex.figs/part2.runnable_directory.3.png}
  \caption{Thin environment shared across multiple runnable directories.}
\end{figure}
% render_images:end

% ----------------------------------------------------------------------------
\subsubsection{Submodule of ``helpers''}

All Causify repositories include a dedicated ``helpers'' repository as
a submodule. This repository contains common utilities and development
toolchains, such as the thin environment, Linter, Docker, and invoke
workflows. By centralizing these resources, we eliminate code
duplication and ensure that all teams, regardless of the project, use
the same tools and procedures.

Additionally, it hosts symbolic link targets for files that must technically
reside in each repository but are identical across all of them (e.g.,
license and certain configuration files). Manually keeping them in
sync can be difficult and error-prone over time. In our approach,
these files are stored exclusively in ``helpers'', and all other repositories
utilize read-only symbolic links pointing to them. This way, we avoid file
duplication and reduce the risk of introducing accidental discrepancies.

% rendered_images:begin
% ```mermaid
% graph RL
%     subgraph A [Runnable Dir A]
%         direction TB
%         B[Runnable Dir B]
%         H1[Helpers]
%     end
%     subgraph H [Helpers]
%     end
% 
%     H -->|Submodule| H1
% 
%     style A fill:#FFF3CD
%     style H fill:#FFF3CD,stroke:#9E9D24
% ```
% caption="Helpers" submodule integrated into a repository.
% rendered_images:end
% render_images:begin
\begin{figure}[H]
  \includegraphics[width=\linewidth]{part2.runnable_directory.tex.figs/part2.runnable_directory.4.png}
  \caption{"Helpers" submodule integrated into a repository.}
\end{figure}
% render_images:end

\paragraph{3.4.1. Git hooks}

Our ``helpers'' submodule includes a set of Git hooks used to enforce policies
across our development process, including Git workflow rules, coding standards,
security and compliance, and other quality checks. These hooks are
installed by default when the user activates the thin environment. They
perform essential checks such as verifying the branch, author
information, file size limits, forbidden words, Python file compilation,
and potential secret leaks\ldots etc.

% ----------------------------------------------------------------------------
\subsubsection{Executing tests}

Our system supports robust testing workflows that leverage the
containerized environment for comprehensive code validation. Tests are
executed inside Docker containers to ensure consistency across development
and production environments, preventing discrepancies caused by
variations in host system configurations. In the case of nested runnable
directories, tests are executed recursively within each directory's
corresponding container, which is automatically identified (see Figure
5). As a result, the entire test suite can be run with a single command,
while still allowing tests in subdirectories to use dependencies that
may not be compatible with the parent directory's environment.

% rendered_images:begin
% ```mermaid
% graph LR
%     start((start))
%     start --> A
%     subgraph A[Runnable Dir A]
%         direction LR
%         pytest_1((pytest))
%         B[Runnable Dir B / Container B]
%         C[Runnable Dir C / Container C]
%         dirA1[dir1 / Container A]
%         dirA2[dir2 / Container A]
%         dirA11[dir1.1 / Container A]
%         dirA12[dir1.2 / Container A]
%         pytest_1 --> B
%         pytest_1 --> C
%         pytest_1 --> dirA1
%         pytest_1 --> dirA2
%         dirA1 --> dirA11
%         dirA1 --> dirA12
%     end
% 
% style A fill:#FFF3CD,stroke:#9E9D24
% style B font-size:15px
% style C font-size:15px
% caption=Recursive test execution in dedicated containers.
% 
% ----------------------------------------------------------------------------
% \subsubsection{Dockerized executables}
% 
% Sometimes, installing a toolchain inside a development container is
% not justified especially when it is large or used only occasionally. In
% such cases, we use \emph{dockerized executables} where the tool and its
% dependencies are packaged into a slim Docker image, and the command is
% executed inside a short-lived container. We launch a short-lived container
% with the repository mounted, run the tool, then remove the container;
% the versioned image stays cached for reuse, avoiding repeated setup.
% 
% This approach prevents the development environment from being bloated with
% rarely used dependencies while preserving reproducibility and version control
% of the toolchain. When needed (e.g., during test execution), a dockerized
% executable can also run inside another Docker container using either the
% children or sibling container approach, as discussed in Section~3.2.
% 
% ============================================================================
% \subsection{Discussion}
% 
% Causify's approach presents a strong alternative to existing code organization
% solutions, offering scalability and efficiency for both small and
% large systems.
% 
% The proposed modular architecture is centered around runnable
% directories, which operate as independent units with their own build and
% release lifecycles. This design bypasses the bottlenecks common in
% large monorepos, where centralized workflows can slow down CI/CD processes
% unless specialized tools like Buck or Bazel are used. By leveraging
% Docker containers, we ensure consistent application behavior across development,
% testing, and production environments, avoiding problems caused by
% system configuration discrepancies. Dependencies are isolated within each
% directory's dedicated container, reducing the risks of issues that
% tight coupling or package incompatibility might create in a monorepo
% or a multi-repo setup.
% 
% Unlike multi-repos, runnable directories can utilize shared utilities from
% ``helper'' submodules, eliminating code duplication and promoting
% consistent workflows across projects. They can even reside under a unified
% repository structure which simplifies codebase management and reduces the
% overhead of maintaining multiple repositories. With support for
% recursive test execution spanning all components, runnable directories
% allow for end-to-end validation of the whole codebase through a single
% command, removing the need for testing each repository separately.
% 
% There are, however, several challenges that might arise in the
% adoption of our approach. Teams that are unfamiliar with containerized
% environments may need time and training to effectively transition to the
% new workflows. The reliance on Docker may introduce additional
% resource demands, particularly when running multiple containers concurrently
% on development machines. This would require further optimization, possibly
% aided by customized tooling. These adjustments, while ultimately beneficial,
% can add complexity to the system's rollout and necessitate ongoing
% maintenance to ensure seamless integration with existing CI/CD
% pipelines and development practices.
% 
% ============================================================================
% \subsection{Future directions}
% 
% Looking ahead, there are several areas where the proposed approach can
% be improved. One direction is the implementation of dependency-aware caching
% to ensure that only the necessary components are rebuilt or retested when
% changes are made. This would reduce the time spent on development
% tasks, making the overall process more efficient. Further optimization
% could involve designing our CI/CD pipelines to execute builds, tests,
% and deployments for multiple runnable directories in parallel, which would
% allow us to take full advantage of available compute resources.
% 
% Additional measures can also be taken to enhance security. Integrating
% automated container image scanning and validation before deployment
% would help guarantee compliance with organizational policies and
% prevent vulnerabilities from entering production environments. In addition,
% fine-grained access controls could be introduced for runnable
% directories in order to safeguard sensitive parts of the codebase. These
% steps will bolster both the security and efficiency of our workflows as
% the projects continue to scale.
% 
% 